% ============================================================
% Chapter 5 — Learning Modes and the RAG Pipeline
% ============================================================
\chapter{Learning Modes and the Retrieval Pipeline}
\label{ch:modes}

\roadmap{This chapter documents the two ways a learner encounters content
and tutoring: Mode A, a free-form dialogue workspace over learner-uploaded
sources; and Mode B, structured lessons with a docked course chatbot. Both
are grounded by a single retrieval-augmented-generation (RAG) pipeline,
which the chapter specifies end to end---chunking, embedding, hybrid
retrieval with rank fusion, optional reranking, and prompt assembly---with
exact constants. It closes with the context-resolution logic that decides
what a tutoring turn is grounded on, including a disclosed inconsistency
between two code paths.}

\section{The two modes}
\label{sec:modes-overview}

The platform exposes a course-level \code{learningMode} of either
\code{STANDARD} or \code{DIALOGUE}.\footnote{\code{Course.learningMode};
\code{apps/api/prisma/schema.prisma}} These correspond to Mode B and Mode
A respectively, and they ground tutoring on different corpora.

\subsection{Mode A: free-form dialogue over learner sources}
\label{sec:modes-a}

Mode A is a NotebookLM-style workspace in which a learner uploads their own
documents and converses with a tutor grounded strictly in those sources.
It comprises a sources panel, a citation-bearing chat panel, a
``studio'' panel that generates study artifacts, and a PDF reader with
highlighting and notes.\footnote{\code{apps/web/src/pages/student/DialogueLearning.tsx}}
Tutoring runs through \code{DialogueService.sendMessage}, which retrieves
from the learner's own chunk corpus (\dbtable{StudentRagChunk}, hard-scoped
to the learner) and persists each turn to
\dbtable{DialogueMessage}.\footnote{\code{apps/api/src/dialogue/dialogue.service.ts}}
The studio generates six artifact types in principle
(\dbtable{StudioOutputType}); five are implemented (briefing document,
flashcard set, table comparison, mind map, FAQ), and the sixth,
\code{TIMELINE}, is a dead enum value with no prompt or schema---disclosed
here and in \cref{ch:limitations}.

\subsection{Mode B: structured lessons with a docked chatbot}
\label{sec:modes-b}

Mode B is the conventional course experience: a teacher authors modules of
items---block-based \code{PAGE}s, \code{PDF} lessons, links, and
assessments---and the learner reads them beside a docked ``Learning
Assistant'' chatbot. The chatbot is the intervention service's
\code{chat()} method; it grounds on the teacher material the learner is
currently viewing---the chunks of the PDF matched by filename, or the text
extracted from a page's block document---and persists to
\dbtable{ChatbotMessage}.\footnote{\code{apps/api/src/learning-interventions/learning-interventions.service.ts}}
Because chatbot conversation is anchored to the learning session but must
outlive it for later review, \dbtable{ChatbotMessage} rows survive a
session purge with a nulled session key (\cref{sec:sensing-coding}).

The distinction that matters analytically is the grounding corpus: Mode A
grounds on \emph{learner-owned} chunks scoped to the individual, Mode B on
the \emph{teacher-owned} course corpus. There is no cross-contamination
between the two.

\section{The retrieval pipeline}
\label{sec:modes-rag}

Both modes share one ingestion-and-retrieval pipeline
(\cref{fig:rag-pipeline}), an instance of the retrieval-augmented-generation
pattern \citep{lewis2020rag}, specified here with the constants taken from
the code.

\begin{figure}[htbp]
\centering
\figplaceholder{RAG pipeline}{Upload $\rightarrow$ chunk (512-token /
100-overlap) $\rightarrow$ embed (per-provider models) $\rightarrow$
hybrid dense+sparse retrieval $\rightarrow$ reciprocal rank fusion
($k=60$) $\rightarrow$ optional Cohere rerank $\rightarrow$ grounded
prompt assembly}
\caption[RAG pipeline]{The shared retrieval-augmented-generation pipeline.
Constants as in \cref{sec:modes-rag}.}
\label{fig:rag-pipeline}
\end{figure}

\subsection{Chunking}
\label{sec:modes-chunking}

A single shared chunker serves both
corpora.\footnote{\code{apps/api/src/rag/shared/chunking.service.ts}} Its
constants are a maximum of 512 tokens per chunk, 100 tokens of overlap, a
four-characters-per-token approximation (so $\approx 2048$ characters per
chunk with $\approx 400$ characters of overlap), and a 50-character
minimum below which a chunk is discarded. The runtime strategy is one of
\code{markdown}, \code{code}, \code{paragraph} (the default), or
\code{auto}, resolved from MIME type or filename. The paragraph strategy
splits on blank lines and re-splits oversize paragraphs by sentence.

A documentation-versus-code discrepancy is disclosed here: the Prisma
enum \code{ChunkingStrategy} still lists \code{FIXED_SIZE},
\code{PARAGRAPH}, and \code{SEMANTIC}, but this enum is vestigial---ingest
no longer consults it, and \emph{no semantic chunking is
implemented}.\footnote{\code{apps/api/src/rag/rag.service.ts:535}}
The report therefore never claims semantic chunking.

\subsection{Embedding}
\label{sec:modes-embedding}

Embedding model and dimensionality are pinned per chunk so that a corpus
can be queried consistently.\footnote{\code{apps/api/src/student-rag/embedding.service.ts};
\code{apps/api/src/llm/model-registry.ts}} The defaults per provider are
OpenAI \code{text-embedding-3-small} (1536 dimensions), Google
\code{gemini-embedding-001} (768), and Cohere \code{embed-v4.0} (1536).
Two optional, off-by-default enrichments exist: \emph{contextual
retrieval}, which prepends a 50--100-token situating blurb to each chunk
before embedding (stored separately and never shown to the generator),
and \emph{multimodal} ingestion, which embeds rendered PDF page images as
Cohere Embed-4 image vectors. A SHA-256 pseudo-embedding fallback exists
for development and is blocked in production.

\subsection{Hybrid retrieval and rank fusion}
\label{sec:modes-retrieval}

Retrieval fuses a dense and a sparse
retriever.\footnote{\code{apps/api/src/student-rag/student-rag-retrieval.service.ts}}
The dense retriever computes cosine similarity in memory over the stored
JSON vectors, grouping chunks by their pinned model/dimension tuple and
re-embedding the query once per group. The sparse retriever is a
lexical/ILIKE scan with a BM25-like score over chunk content, contextual
blurb, and image caption. The two ranked lists are combined by reciprocal
rank fusion with the constant $k=60$ \citep{cormack2009rrf}: for a chunk
at rank $r$ in a list, its fused contribution is
\begin{equation}
  \text{score} = w \cdot \frac{1}{k + r + 1}, \qquad k = 60,
\end{equation}
with $w$ a modality weight (unity by default). Near-duplicate chunks
(word-level Jaccard overlap above $0.9$) are dropped.

An optional Cohere reranking stage (model \code{rerank-english-v3.5}) can
reorder the fused candidates; it is \emph{off by default}, and its no-op
fallback preserves the fusion order whenever the reranker is disabled,
unkeyed, or times out. A documented inconsistency is disclosed: schema
fields named \code{rerankTopK} exist on both the course and the dialogue
settings but are \emph{not read anywhere}; the effective top-$k$ is a
literal at each call site (10 for teacher retrieval, 5 for interventions,
8 for dialogue).\footnote{Verified by tree-wide search;
\code{apps/api/src/rag/rag.service.ts}, \code{dialogue.service.ts}.}

\subsection{Prompt assembly and persistence}
\label{sec:modes-prompt}

A single grounding contract assembles every tutoring
prompt.\footnote{\code{apps/api/src/rag/shared/grounded-prompt.ts}} The
system message pairs a persona with fixed grounding rules---answer only
from supplied sources, cite each claim with a source label in a fixed
format, and state explicitly when the sources do not contain the
answer---and the context is a labeled block of retrieved chunks, with
image chunks attached as data URLs where multimodal is enabled.
Dialogue turns persist to \dbtable{DialogueMessage} within a transaction
(carrying parsed citations and token usage); chatbot turns persist to
\dbtable{ChatbotMessage} fire-and-forget (carrying the context source,
selected text, current page, and any suggested strategy).

\section{Context resolution and a disclosed inconsistency}
\label{sec:modes-context}

What a tutoring turn is grounded on is decided by a precedence chain.
For the interventions, \code{resolveInterventionContext()} resolves in the
order \emph{selection} $\rightarrow$ \emph{PDF source} $\rightarrow$
\emph{student RAG}, and throws if nothing
resolves.\footnote{\code{learning-interventions.service.ts:760}} A text
selection is accepted as the grounding source only when it is at least 20
characters after boilerplate stripping---the \emph{$\geq$20-character
selection rule} referenced throughout this report.

The docked chatbot's \code{chat()} method uses the \emph{same}
20-character rule but a \emph{different} precedence: \emph{PDF source}
(when a content id is present) $\rightarrow$ \emph{selection} $\rightarrow$
\emph{student RAG}, recording a context source of \code{none} rather than
throwing when nothing
resolves.\footnote{\code{learning-interventions.service.ts:4027,4178}}
The consequence, disclosed plainly: when a learner has both an active
selection and an open PDF, an intervention grounds on the \emph{selection}
while the chatbot grounds on the \emph{PDF}. This divergence is recorded
in the per-turn \code{contextSource} field, so its effect on any analysis
of grounding is itself measurable.

\medskip
\noindent The retrieval substrate in place, the next chapter turns to the
four strategy interventions that sit on top of it.
